%=============================================================================
% EXPERIMENTAL SETUP
%=============================================================================
\section{Experimental Setup}
\label{sec:setup}

\subsection{Dataset overview}
The corrected series (\S\ref{sec:corrected}) covers 25 districts over 559 weeks,
2013-W26 to 2024-W10, with weekly cases, ERA5 temperature, precipitation, humidity
and soil moisture, population, and the benchmark's district graph. Seven weeks are
missing and windows that touch them are dropped. The extension (\S\ref{sec:newweeks})
adds 127 weeks to 2026-W32 for the prospective test only. 
Every model, baseline and control is scored on exactly the same validation and test
windows and trained on the same training windows, except that arms with climate lags
skip the first training weeks, which have no climate history that far back.

\subsection{Baselines}
\label{sec:baselines}
Following the brief's baseline rule (LSTM for time series, GCN for graphs), and the
literature on this dataset, we compare against:
\begin{itemize}\itemsep1pt
\item \textbf{Persistence} $\hat y_{t+h}=y_t$ and \textbf{seasonal naive}
      $\hat y_{t+h}=y_{t+h-52}$, the classical floors.
\item \textbf{LSTM}, the two-layer encoder of SEIR-LSTM~\cite{liu2025seirlstm}, and
      a \textbf{GCN}~\cite{kipf2017gcn} against the same network without message
      passing, the mandatory sequence and graph baselines.
\item The \textbf{five published ST-GNNs} of the Sri Lankan benchmark~\cite{weng2024graph}:
      STGAT, A3TGCN, ASTGCN, AAGCN and DCRNN, each with its direct head as published.
\item \textbf{SEIR-informed models}: the SEIR-LSTM~\cite{liu2025seirlstm}, its graph
      version (SEIR-GNN), a pure SEIR decoder, and a metapopulation SEIR head in the
      style of MepoGNN~\cite{cao2023mepognn} and CausalGNN~\cite{wang2022causalgnn}.
\item \textbf{Classical and simple learners}: gradient-boosted trees with and without
      climate~\cite{chen2016xgboost}, a negative-binomial GLM, $k$-nearest-neighbour
      analogues, an AR(3) ridge regression, and an STID-style MLP~\cite{shao2022stid}.
\end{itemize}

\subsection{Evaluation protocol}
\label{sec:protocol}
\textbf{Splits.} Rolling origins at 0.55, 0.70 and 0.85 of the series, three seeds each.
The 30 windows before each origin are validation, the following 15\% are test, and
normalisation uses training weeks only; no window is shuffled across time. A
confirmation uses nine disjoint origins at $0.40+k/15$ ($k=0,\dots,8$). The
prospective test trains once on all development weeks and scores the 85 forecast
starts from 2024-W11 onward, with no target week shared between partitions.

\textbf{Metrics.} RMSE in weekly cases, pooled over every (window, district, horizon)
cell, is primary; MAE, SMAPE and MAPE (on targets $\ge1$) are secondary, overall and
per horizon. Tables report the mean over runs.

\textbf{Selection and tests.} Selection uses validation RMSE only; test numbers are
written to disk and never used to choose. The protocol and, for each experiment, its
arms, control, decision rule and predicted outcome were committed before it ran. An
arm is adopted if its mean validation difference to its control is negative, it wins
at least 7 of 9 (origin, seed) units, and it does not raise horizon-3 error.
Differences are paired on matched (origin, seed) units and tested with an exact
two-sided sign-flip permutation test, BH-adjusted within each family. With three
origins this unit mostly measures run-to-run stability; claims about the series need
the nine-origin protocol, where the origin-level test can reach $p=0.004$. The
prospective test uses a block bootstrap over forecast starts (block 8, 10{,}000
replicates) with Holm correction over its two pre-registered comparisons.
